\section{Theoretical Adaptation}
\label{sec:theoretical_adaptation}

This chapter presents three proposals for adapting inconsistency measures to classical planning problems. This work develops planning-native measures that operate directly on planning structures: goals, operators, and precondition-effect relationships. The approach employs multi-faceted diagnostic profiles rather than single severity scores, providing actionable information for problem repair without making assumptions about which repairs are valid.

The chapter begins with six test scenarios (Section~\ref{sec:test_scenarios}) that serve as running examples throughout the development, followed by motivation for planning-native measures (Section~\ref{sec:naive_encoding}). Sections~\ref{sec:unreachability_profile}--\ref{sec:sequencing_conflict_profile} present the proposals informally, with formal definitions in Section~\ref{sec:formalization}, rationality postulates in Section~\ref{sec:postulates}, and complexity analysis in Section~\ref{sec:complexity}.

\subsection{Test Scenarios}
\label{sec:test_scenarios}

This section introduces six test scenarios that serve as running examples throughout the theoretical development. Each scenario represents a minimal planning problem exhibiting a specific structural pattern of unsolvability. The scenarios are organized by the measure designed to detect their failure mode, though some scenarios (particularly sequencing conflicts) trigger multiple measures.

All scenarios are expressed as STRIPS planning problems $\langle P, O, I, G \rangle$ following the notation established in Section~\ref{sec:foundations_planning}. For each scenario, the intuitive description, formal specification, and structural pattern are provided.

\subsubsection{Unreachability Scenarios}

The following scenarios exhibit dependency structure failures where required propositions never appear in any reachable state. They differ in dependency depth (chain length) versus width (parallel chains).

\paragraph{Scenario 1: Locked Door.}

An agent must enter a room, but entry requires unlocking a door, which requires a key that no operator produces. This creates a linear dependency chain of depth~2.

\begin{align*}
  P & = \{\text{has\_key}, \text{door\_unlocked}, \text{in\_room}\}                            \\
  O & \ni \text{unlock} = \langle \{\text{has\_key}\}, \{\text{door\_unlocked}\}, \{\} \rangle \\
    & \ni \text{enter} = \langle \{\text{door\_unlocked}\}, \{\text{in\_room}\}, \{\} \rangle  \\
  I & = \{\}                                                                                   \\
  G & = \{\text{in\_room}\}
\end{align*}

The dependency chain $\textit{in\_room} \leftarrow \textit{door\_unlocked} \leftarrow \textit{has\_key}$ fails at the root: no operator produces \textit{has\_key}. Forward reachability yields $S_{\text{reach}} = \{\{\}\}$.

\paragraph{Scenario 2: Bank Vault.}

A more complex variant requires three independent precondition chains (two keys and admin approval) to achieve a single goal. This creates parallel dependency failures with width~3.

\begin{align*}
  P & = \{\text{key1}, \text{key2}, \text{admin\_approval},                                                                             \\
    & \quad \text{door1\_open}, \text{door2\_open}, \text{access\_granted}, \text{in\_vault}\}                                          \\
  O & \ni \text{open\_door1} = \langle \{\text{key1}\}, \{\text{door1\_open}\}, \{\} \rangle                                            \\
    & \ni \text{open\_door2} = \langle \{\text{key2}\}, \{\text{door2\_open}\}, \{\} \rangle                                            \\
    & \ni \text{grant\_access} = \langle \{\text{admin\_approval}\}, \{\text{access\_granted}\}, \{\} \rangle                           \\
    & \ni \text{enter} = \langle \{\text{door1\_open}, \text{door2\_open}, \text{access\_granted}\}, \{\text{in\_vault}\}, \{\} \rangle \\
  I & = \{\}                                                                                                                            \\
  G & = \{\text{in\_vault}\}
\end{align*}

Three independent chains fail at their roots (\textit{key1}, \textit{key2}, \textit{admin\_approval} are all missing producers), blocking all downstream propositions. Forward reachability yields $S_{\text{reach}} = \{\{\}\}$.

\subsubsection{Mutex Scenarios}

The following scenarios exhibit operator-induced mutual exclusions where goal propositions are each individually achievable but never simultaneously true in any reachable state. Actions are reversible: the system can transition between goal-satisfying states, but effects structurally prevent co-occurrence.

\paragraph{Scenario 3: Light Switch.}

A goal requires both \text{on} and \text{off} to be true simultaneously, but toggle actions enforce mutual exclusion.

\begin{align*}
  P & = \{\text{on}, \text{off}\}                                                         \\
  O & \ni \text{turn\_on} = \langle \{\text{off}\}, \{\text{on}\}, \{\text{off}\} \rangle \\
    & \ni \text{turn\_off} = \langle \{\text{on}\}, \{\text{off}\}, \{\text{on}\} \rangle \\
  I & = \{\text{off}\}                                                                    \\
  G & = \{\text{on}, \text{off}\}
\end{align*}

Forward reachability yields $S_{\text{reach}} = \{\{\text{off}\}, \{\text{on}\}\}$. Both goals appear in reachable states, but no state contains both. Actions are reversible (can toggle indefinitely), making this a pure mutex without sequencing conflict.

\paragraph{Scenario 4: Traffic Light.}

A goal requires all three colors (\textit{red}, \textit{yellow}, \textit{green}) simultaneously, but cyclic transitions enforce that exactly one color is active at any time.

\begin{align*}
  P & = \{\text{red}, \text{yellow}, \text{green}\}                                                        \\
  O & \ni \text{red\_to\_green} = \langle \{\text{red}\}, \{\text{green}\}, \{\text{red}\} \rangle         \\
    & \ni \text{green\_to\_yellow} = \langle \{\text{green}\}, \{\text{yellow}\}, \{\text{green}\} \rangle \\
    & \ni \text{yellow\_to\_red} = \langle \{\text{yellow}\}, \{\text{red}\}, \{\text{yellow}\} \rangle    \\
  I & = \{\text{red}\}                                                                                     \\
  G & = \{\text{red}, \text{yellow}, \text{green}\}
\end{align*}

Forward reachability yields $S_{\text{reach}} = \{\{\text{red}\}, \{\text{green}\}, \{\text{yellow}\}\}$. All three goals are individually achievable (the cycle is complete), but no state contains more than one. This forms a complete mutex clique: every pair of goals is mutually exclusive.

\subsubsection{Sequencing Conflict Scenarios}

The following scenarios exhibit irreversible exclusions where achieving certain goal propositions creates permanent dead-end states. Unlike mutex scenarios where goals alternate via reversible actions, sequencing conflicts involve one-way transitions. Consistent with the measure hierarchy (Proposition~\ref{prop:measure_hierarchy}), these scenarios exhibit both mutex and sequencing conflict values.

\paragraph{Scenario 5: Trust and Travel.}

An agent must secure a local partnership that requires established trust, and exploit a remote opportunity. Traveling destroys the trust relationship irreversibly. This creates a unidirectional sequencing conflict.

\begin{align*}
  P & = \{\text{at\_local}, \text{at\_remote}, \text{trust}, \text{partnership}, \text{opportunity}\}                                         \\
  O & \ni \text{secure} = \langle \{\text{at\_local}, \text{trust}\}, \{\text{partnership}\}, \{\} \rangle                                    \\
    & \ni \text{travel} = \langle \{\text{at\_local}\}, \{\text{at\_remote}\}, \{\text{at\_local}, \text{trust}, \text{partnership}\} \rangle \\
    & \ni \text{exploit} = \langle \{\text{at\_remote}\}, \{\text{opportunity}\}, \{\} \rangle                                                \\
    & \ni \text{return} = \langle \{\text{at\_remote}\}, \{\text{at\_local}\}, \{\text{at\_remote}\} \rangle                                  \\
  I & = \{\text{at\_local}, \text{trust}\}                                                                                                    \\
  G & = \{\text{partnership}, \text{opportunity}\}
\end{align*}

Both goals are individually achievable: \textit{partnership} via \textit{secure}, and \textit{opportunity} via \textit{travel} $\rightarrow$ \textit{exploit}. However, the goals never coexist (mutex). The conflict is asymmetric:
\begin{itemize}
  \item $(\text{opportunity}, \text{partnership})$ \textbf{is} a sequencing conflict: from any state with \textit{opportunity}, the agent lacks \textit{trust} to achieve \textit{partnership} (travel destroyed it).
  \item $(\text{partnership}, \text{opportunity})$ is \textbf{not} a sequencing conflict: from a \textit{partnership}-state, \textit{opportunity} remains reachable via \textit{travel} $\rightarrow$ \textit{exploit} (though this path deletes \textit{partnership}).
\end{itemize}

\paragraph{Scenario 6: Rival Alliances.}

An agent seeks alliances with two rival factions, but each faction demands exclusive loyalty. Allying with one faction alienates the other. This creates a symmetric sequencing conflict.

\begin{align*}
  P & = \{\text{neutral\_A}, \text{neutral\_B}, \text{allied\_A}, \text{allied\_B}\}                          \\
  O & \ni \text{ally\_A} = \langle \{\text{neutral\_A}\}, \{\text{allied\_A}\}, \{\text{neutral\_B}\} \rangle \\
    & \ni \text{ally\_B} = \langle \{\text{neutral\_B}\}, \{\text{allied\_B}\}, \{\text{neutral\_A}\} \rangle \\
  I & = \{\text{neutral\_A}, \text{neutral\_B}\}                                                              \\
  G & = \{\text{allied\_A}, \text{allied\_B}\}
\end{align*}

Forward reachability yields $S_{\text{reach}} = \{\{\text{neutral\_A}, \text{neutral\_B}\}, \{\text{neutral\_A}, \text{allied\_B}\}, \{\text{neutral\_B}, \text{allied\_A}\}\}$. Both goals are individually achievable but never coexist (mutex). Unlike Trust and Travel, both directions are sequencing conflicts: allying with either faction alienates the other, deleting the precondition for the remaining alliance.

\subsubsection{Scenario Summary}

Table~\ref{tab:scenario_summary} summarizes the test scenarios and their structural characteristics.

\begin{table}[ht]
  \centering
  \caption{Test scenarios and their structural patterns}
  \label{tab:scenario_summary}
  \small
  \begin{tabular}{lllll}
    \hline
    Scenario         & Category & Pattern         & Measure         & Notes         \\
    \hline
    Locked Door      & 2a       & Linear chain    & $I_{\text{UR}}$ & Depth 2       \\
    Bank Vault       & 2a       & Parallel chains & $I_{\text{UR}}$ & Width 3       \\
    Light Switch     & 2c       & Simple pair     & $I_{\text{MX}}$ & Reversible    \\
    Traffic Light    & 2c       & Complete clique & $I_{\text{MX}}$ & 3 mutex pairs \\
    Trust and Travel & 2c       & Unidirectional  & $I_{\text{GS}}$ & One-way block \\
    Rival Alliances  & 2c       & Symmetric       & $I_{\text{GS}}$ & Mutual block  \\
    \hline
  \end{tabular}
\end{table}

These scenarios will be referenced throughout the remainder of this chapter to illustrate measure definitions (Section~\ref{sec:proposed_measures}), demonstrate postulate compliance (Section~\ref{sec:postulates}), and provide concrete measure values for comparison (Table~\ref{tab:measure_values}).


\subsection{Motivation: Limitations of Direct Measure Application}
\label{sec:naive_encoding}

A natural approach to measuring inconsistency in planning problems would encode the planning problem as a propositional knowledge base and apply classical inconsistency measures. Such an encoding would include: (1)~the goal conjunction $\bigwedge_{g \in G} g$, (2)~operator-derived mutex constraints, (3)~precondition dependencies as implications, and (4)~closed-world assertions $\neg p$ for propositions not in $I$. For the Locked Door scenario (Section~\ref{sec:test_scenarios}), this yields implications $\textit{in\_room} \rightarrow \textit{door\_unlocked} \rightarrow \textit{has\_key}$ combined with $\neg\textit{has\_key}$, producing a conflict resolvable by assigning $B$ (``both true and false'') to $\textit{has\_key}$, hence $I_c = 1$.

However, this naive encoding yields uninformative results. Table~\ref{tab:naive_ic} shows that $I_c = 1$ for representative scenarios from each structural category.

\begin{table}[ht]
    \centering
    \caption{Contension measure $I_c$ for naive propositional encodings}
    \label{tab:naive_ic}
    \begin{tabular}{llc}
        \hline
        Category                & Scenario                      & $I_c$ \\
        \hline
        2a: Missing Producers   & Locked Door (Scenario 1)      & 1     \\
        2c: Mutex               & Light Switch (Scenario 3)     & 1     \\
        2c: Sequencing Conflict & Trust and Travel (Scenario 5) & 1     \\
        \hline
    \end{tabular}
\end{table}

This uniform result fails to distinguish fundamentally different types of unsolvability. The cause lies in three-valued semantics: the mutual exclusivity constraint $\neg(p \land q)$ is satisfied by assigning the conflicting value $B$ to only \emph{one} proposition, since $\neg(B \land T) = \neg B = B$. Thus, symmetric conflicts (where both propositions are equally problematic) yield the same $I_c$ value as single-point failures.

This finding motivates the development of \emph{planning-native} inconsistency measures that operate directly on planning structures (goals, operators, reachable states) rather than reducing to propositional logic. The measures proposed in this chapter follow the inconsistency measurement methodology (formal definitions, rationality postulates, complexity analysis) while respecting planning-specific semantics.

\subsection{Design Philosophy and Approach}

The adaptation strategy is guided by several key principles that distinguish this work from traditional inconsistency measurement approaches. Rather than reducing inconsistency to a single numerical severity score, the proposed measures provide \emph{diagnostic profiles} consisting of multiple complementary values. This design choice recognizes that different aspects of unsolvability (such as the number of affected goals versus the extent of broken dependencies) serve different debugging purposes.

A fundamental principle guiding these proposals is to ``measure what IS, not what SHOULD BE.'' Traditional repair-based inconsistency measures implicitly encode assumptions about valid repairs by counting minimal changes needed to restore consistency. In planning contexts, however, determining which components ``should'' exist (missing operators, initial state facts, or goal modifications) requires domain expertise that automated measures cannot provide. The proposed measures therefore report objective structural properties (reachable propositions, action-based exclusions) without presupposing which repairs are appropriate.

To maximize utility for planning practitioners, the measures employ planning-specific concepts (goals, operators, reachability, preconditions) rather than generic logical terminology (formulas, interpretations, models). The structural taxonomy established in \ref{sec:taxonomy} suggests reachability analysis for dependency structures (Category~2).

\subsection{From Structural Categories to Measures}

The structural taxonomy organizes unsolvability causes by their mathematical properties: goal specification conflicts (Category~1), dependency graph failures (Category~2), and resource insufficiency (Category~3). The measurement framework below focuses on Category~2 problems with three measures forming a diagnostic hierarchy.

The measures differ in their analysis scope. The first two measures (\emph{Unreachability} and \emph{Mutex}) perform \textbf{state space analysis}: they compute reachability once and analyze properties of the resulting state set, asking ``do the required states exist?'' The third measure (\emph{Sequencing}) performs \textbf{trajectory analysis}: it examines path-dependent properties from goal-achieving states, asking ``can we reach other goals from here?''

Together, these three measures provide complementary diagnostic perspectives on Category~2 dependency problems, with a key relationship: sequencing conflicts imply mutex (achieving one goal blocks another implies they never coexist), so the measures form a hierarchy rather than detecting independent phenomena.

\subsection{Measurement Framework}

Each proposal defines a \emph{diagnostic profile} consisting of two complementary \emph{measures}. Following terminology in the inconsistency measurement literature~\cite{Thimm2019}, a measure assigns a single numerical value to a knowledge base (here, a planning problem). Each measure provides two values: one indicating \emph{scope} (how many goals are affected) and one revealing \emph{structure} (underlying complexity). This design provides richer diagnostic information than single values while maintaining interpretability.

The three proposals are:

\emph{Proposal 1: Unreachability Profile} reports propositions that never appear in any reachable state, providing a generic diagnostic for dependency structure problems. The profile reports unreachable goal propositions (indicating which goals are blocked) and total unreachable propositions (revealing dependency depth). This measure detects unreachability regardless of structural cause (missing producers, circular dependencies, or action-based exclusions) but does not distinguish between these causes.

\emph{Proposal 2: Goal Mutex Profile} detects operator-induced mutual exclusions where goal propositions are each individually achievable but never simultaneously true in any reachable state. The profile reports which goal propositions participate in mutex relationships and counts unordered mutex pairs, revealing fundamentally incompatible goals even when actions allow transitions between them.

\emph{Proposal 3: Goal Sequencing Conflict Profile} detects irreversible dead-end states where achieving certain goal propositions permanently blocks others. Unlike mutex (where goals exclude each other but actions are reversible), sequencing conflicts involve one-way transitions that prevent subsequent goal achievement. The profile reports which goal propositions participate in conflicts and counts ordered conflict pairs, revealing directionality. Since blocking implies exclusion, sequencing conflicts are a strict subset of mutex relationships.

\paragraph{Scope and Structure Variants.}

Each measure provides two complementary values following a standardized pattern: the \emph{scope} variant reports goal-level severity, while the \emph{structure} variant reveals underlying complexity. This standardization enables consistent interpretation across measures while allowing each structure variant to capture the natural ``next level of detail'' for its specific failure mode.

The scope values are directly comparable across measures: a problem with 3 unreachable goals and 2 mutex goals immediately indicates which failure mode is more prevalent. The structure values are not directly comparable but provide diagnostic depth within each measure: the ratio between scope and structure reveals characteristic patterns (dependency chains, conflict hubs, directional asymmetries).

Specifically, structure variants differ by analysis type:
\begin{itemize}
    \item \textbf{Unreachability}: Structure expands \emph{vertically} from goals to all propositions, revealing dependency depth
    \item \textbf{Mutex}: Structure expands \emph{horizontally} to count pairwise relationships, revealing conflict complexity
    \item \textbf{Sequencing}: Structure counts \emph{ordered} pairs, revealing directional conflict patterns
\end{itemize}

This variation is a feature, not an inconsistency. Each structure variant answers ``what is the broader context of this failure mode?'' in the most informative way for that specific analysis.

\subsection{Proposal 1: Unreachability Profile}
\label{sec:unreachability_profile}

The first proposal provides a unified approach to Category~2 unsolvability by characterizing unreachable propositions through reachability analysis. This measure applies uniformly to all dependency structure subtypes: missing producers (2a), circular dependencies (2b), and dead-end states (2c).

The \emph{unreachability profile} employs forward reachability analysis to compute: (1) the number of goal propositions not appearing in any reachable state, and (2) the total number of propositions not appearing in any reachable state. The first value identifies which goal propositions are blocked, while the second reveals the scope of broken dependencies. The difference between these values indicates cascading unreachability depth.

The Locked Door scenario (Scenario~1) yields 1 unreachable goal but 3 total unreachable propositions, revealing a dependency chain of depth~2. The Bank Vault scenario (Scenario~2) yields 1 unreachable goal but 7 total unreachable propositions, revealing three parallel missing-producer failures.

A small difference between scope and structure indicates shallow failures; a large difference reveals deep cascading dependencies. This measure applies uniformly across all Category~2 subtypes, providing a generic diagnostic for dependency-related unsolvability.

\subsection{Proposal 2: Goal Mutex Profile}
\label{sec:goal_mutex_profile}

The second proposal addresses unsolvability where goal propositions are each individually achievable but never simultaneously true in any reachable state. This differs from specification-level contradictions (goals containing explicit negation like $\{p, \neg p\}$, which do not exist in propositional STRIPS) and from sequencing conflicts (where achieving one goal permanently blocks another). Mutex relationships are \emph{operator-induced}: actions enforce mutual exclusion through their effects, creating states where goals alternate but never coexist.

The \emph{goal mutex profile} characterizes these operator-induced exclusions through reachability analysis. Specifically, goal propositions $g_1$ and $g_2$ are \emph{mutex} when: (1) $g_1$ appears in some reachable state, (2) $g_2$ appears in some reachable state, but (3) no reachable state contains both $g_1$ and $g_2$ simultaneously. This captures situations where goals are individually achievable and actions may allow transitions between them (reversibility), but the action model structurally prevents their co-occurrence.

The profile computes two values: (1) the number of goal propositions participating in at least one mutex relationship, and (2) the total number of unordered mutex pairs. The first value indicates scope, while the second reveals conflict structure (isolated pairs versus complete conflict cliques).

The Light Switch scenario (Scenario~3) yields 2 mutex goals and 1 mutex pair: both propositions are individually achievable and actions are reversible, but no state contains both. The Traffic Light scenario (Scenario~4) yields 3 mutex goals and 3 mutex pairs (a complete clique): all three colors are individually achievable via cyclic transitions, but exactly one is active at any time.

This measure identifies operator-induced mutual exclusions where goals cannot coexist despite each being individually achievable. The ratio of mutex pairs to mutex goals distinguishes isolated conflicts from dense cliques.

\subsection{Proposal 3: Goal Sequencing Conflict Profile}
\label{sec:sequencing_conflict_profile}

The third proposal addresses unsolvability of Category~2c: sequencing conflicts where achieving certain goal propositions creates dead-end states that prevent other goal propositions from being reached. Unlike mutex relationships (Proposal~2) where goals can be achieved in different states but never simultaneously, sequencing conflicts involve \emph{irreversible} state transitions where achieving one goal permanently blocks another.

The \emph{goal sequencing conflict profile} characterizes action-based exclusions through reachability analysis. Specifically, a \emph{sequencing conflict} $(g_1, g_2)$ exists when: (1) both goal propositions $g_1$ and $g_2$ can each be achieved from the initial state, but (2) from any reachable state where $g_1$ is true, goal proposition $g_2$ is no longer reachable. This captures situations where achieving $g_1$ through available actions inevitably transitions to states from which $g_2$ cannot be satisfied, typically because actions required to achieve $g_1$ delete preconditions needed for $g_2$ without any operator to restore them. The requirement that both goals be achievable distinguishes meaningful sequencing conflicts from cases where $g_2$ was never achievable.

The profile computes two values: (1) the number of goal propositions participating in at least one sequencing conflict (as either blocker or blocked), and (2) the total number of ordered conflict pairs $(g_1, g_2)$. The first value indicates scope, while the second reveals conflict patterns (star structures versus symmetric exclusions).

The Trust and Travel scenario (Scenario~5) yields 2 conflicting goals and 1 ordered conflict pair: exploiting the remote opportunity permanently blocks securing the partnership, but not vice versa. The Rival Alliances scenario (Scenario~6) yields 2 conflicting goals and 2 ordered conflict pairs (both directions), distinguishing symmetric from asymmetric irreversibility.

This measure provides directional information beyond mutex. Sequencing conflicts are a subset of mutex relationships (blocking implies non-coexistence), but the ordered pair count reveals \emph{which} goal achievement causes the dead-end. Reversible mutex problems (Light Switch, Scenario~3) exhibit mutex without sequencing; irreversible problems (Trust and Travel, Scenario~5) exhibit both.

\subsection{Normalization and Relative Measures}

The three proposals presented above report \emph{absolute} inconsistency values: counts of unreachable propositions, mutex pairs, and sequencing conflict pairs. An important consideration for practical deployment is whether these measures should be \emph{normalized} or made \emph{relative} to problem characteristics such as the total number of goal propositions or the size of the proposition space.

The literature on relative inconsistency measures~\cite{Besnard2020} explores this question in the context of propositional knowledge bases, demonstrating both advantages and limitations of normalization approaches. Relative measures can facilitate comparison across problems of different sizes, potentially providing a notion of ``percentage inconsistency'' analogous to test coverage metrics in software engineering.

However, several considerations suggest that absolute measures may be more appropriate for planning diagnostics. First, the severity of unsolvability is not necessarily proportional to problem size: a planning problem with 100 goal propositions where only 2 conflict may be easier to repair than a problem with 3 goal propositions where all 3 form a complete conflict graph. Second, normalization requires choosing appropriate denominators (total goal propositions? total propositions in the problem? maximum possible conflicts?), and different choices encode different assumptions about problem structure. Third, the proposed measures already provide multiple values that capture relative information: the ratio between unreachable goal propositions and total unreachable propositions indicates dependency depth, while the ratio between mutex goals and mutex pairs (or conflicting goals and sequencing pairs) reveals conflict structure.

For the scope of this thesis, the focus remains on absolute measures that report objective structural properties. Future work may explore normalization schemes tailored to specific use cases, such as comparing inconsistency trends across problem generator parameters or establishing thresholds for automated problem classification. The formalization in the following section defines the measures in their absolute form.
